\documentclass[a4paper,10pt]{article}

\usepackage[utf8]{inputenc}
\usepackage[T1]{fontenc}
\usepackage{lmodern}
\usepackage{amsthm}
\usepackage{amsmath}
\usepackage[top=2cm, bottom=2cm, left=2cm, right=2cm]{geometry}
\usepackage{fancyhdr}
\usepackage{draftwatermark}
\usepackage{eurosym}
\usepackage{graphicx}
\usepackage[svgnames]{xcolor}
\usepackage{listings}
\usepackage[francais]{babel}
\pagestyle{fancy}

\SetWatermarkText{\includegraphics{../../logo.png}}
\SetWatermarkScale{0.5}
\SetWatermarkAngle{0}


%opening
\title{\textbf{Cours de géopolitique}}
\date{\today}

\fancyhf{}
\renewcommand{\headrulewidth}{0pt}
\fancyfoot[C]{\tiny Che Bedara - BDE Télécom ParisTech}
\fancyfoot[RO]{\thepage}
\fancyfoot[LE]{\thepage}


\begin{document}

\maketitle

\tableofcontents

\newpage

\part{De quoi s'agit-il ?}

\section{L'invention d'une nouvelle discipline scientifique}

	\subsection{Les fondateurs}

		Le mot géopolitique, depuis son invention, la dernière année du XIX\up{e} siècle, par le professeur suédois de science politique, Rudolf Kjellen
	\footnote{Rudolf Kjellen, \emph{Stormakterna (les Grandes Puissances)}, Stockholm, 1905 ; \emph{Staten som livsform (L’État comme organisme vivant)}, Stockholm, 1920.}
	(1864-1922), a connu, selon les lieux et les époques, des fortunes diverses, liées au sens qui lui a été donné et à l’emploi qui en a été fait.

		Les travaux de Johan Rudolf Klellén sont inspirés par le géographe allemand Friedrich Ratzel, nationaliste, pour le pan-germanisme.
		Friedrich Ratzel est l'un des fondateurs de la géographie politique.
		Il est l'auteur de l'ouvrage « Erde und das Leben », avec le concept de « Lebensraum » qui va dans le sens de l'expansionnisme et sera exploité plus tard par les nazis.

		À la même époque on trouve également des penseurs comme Halford Mackinder (1861-1947) : « Who rules East Europe commands the Heartland ; who rules the Heartland commands the World Island ; who rules the World Island commands the World ».
		Il attache surtout de l'importance au contrôle de la terre et des ressources (agricole, énergétique,...).

		Une autre réflexion intéressante est celle de Alfred Thayer Mahan (américain, 1840-1914).
		En tant qu'ancien amiral il attache lui une importance majeure au contrôle des voies maritimes.
		Il cherche en particulier à expliquer le succès de l'État anglais malgré une population et des ressources modestes.

	\subsection{Le fruit du contexte « fin de siècle »}

		\paragraph{Premier facteur :}
		L'émergence de la géopolitique se fait dans le contexte de scientisme : une foi très forte dans le progrès, qui pousse même certains à croire être arrivés presque à la fin des découvertes scientifiques possibles.

		Autre point important : le darwinisme social.
		Le darwinisme naturel amène l'idée que les plus forts vont éliminer les plus faibles, participant ainsi à l'évolution.
		Dans l'ébullition scientifique de l'époque, il est alors tentant d'appliquer aux groupes sociaux les principes de la sélection naturelle : les plus organisés et les plus dynamiques qui seraient les plus forts seraient appelés à survivre aux autres.

		\paragraph{Second facteur :}
		La perception des distances terrestres est également relativisée. \emph{Le tour du monde en 80 jours} de Jules Verne paraît en 1873.
		Des progrès sont encore faits à la fin du XIX\up{e} siècle en matière de transports et de communications (télégraphe).

		Des questions peuvent alors se poser à l'échelle planétaire.

		\paragraph{Troisième facteur :}
		Le développement des nationalistes.

		L'allemand Herder, à rebours des nationalistes français, explique que la nation est une question de naissance et non de choix. La nation de quelqu'un s'impose à lui et n'est pas choisie.

		Il est plus souvent admis qu'on ne peut pas appartenir à deux nations à la fois.
		L'idée du nationalisme émerge alors.
		Exemples : ligue des patriotes de Paul Déroulède en France, pan-germanisme en Allemagne.

		L'empereur Guillaume II d'Allemagne utilise notamment l'expression “Weltpolitik” : la politique se pense désormais à l'échelle du monde.
		On a également des appétits coloniaux : l'Allemagne ne parvient pas à coloniser le Maroc face aux français et aux anglais.
		Enfin à l'époque un jeu d'alliances se mets en place, que l'on retrouvera lors de la grande guerre.

	\subsection{Une géopolitique plus politique que géographique ?}

		\paragraph{Les conséqences}

		La géographie se met alors au service des ambitions nationalistes : cartes prévisionelles de conquêtes en Europe, découpage coloniale.

		Après la guerre la géopolitique est dévoyée par Karl Haushofer.
		Des liens toxiques s'établissent entre ses idées et le parti nazi.
		Son influence s'exerce surtout à travers le journal \emph{Geopolitik}, très lu parmi ses confrères.
		Il s'en sert de tribune pour s'exprimer contre le traité de Versailles.

		Jacques Ancel, géographe français ayant le plus écrit sur la géopolitique, avec notamment les ouvrages \emph{Géographie des frontières} et \emph{Géopolitique}, dénonce à quel point les travaux de Haushofer sont biaisés.
		Les thèses de Haushofer suscitent donc un vif débat.

		Loin de ces réflexions européennes, l'américain Nicholas Spykman repart des thèses de Mackinder, avec un découpage proche, mais en concentrant son attention sur le contrôle du premier croissant.

		\paragraph{Une discipline discréditée ?}

		On commence à parler de la géopolitique comme pseudo-science.
		Un colloque à la Sorbonne conclut même que l'on devrait cesser d'utiliser ce terme et de l'enseigner.

		Toutefois les pratiques de la géopolitique demeurent : dans l'exploitation du milieu par des conquérants notamment.
		Exemples : bataille de Salamine (victoire de Thémistocle par la connaissance de la géographie), traversée des Alpes par Hannibal avec des éléphants, victoire de Napoléon à Austerlitz (qui amène la cavalerie ennemie sur des lacs gelés, avec saisie des canons aux ennemis qui serviront à construire la colonne Vendôme).

		Les grands états ont toujours des préoccupations géopolitiques après 1945.
		La pratique précède et survit au concept.

		En 1979 a lieu l'invasion du Cambodge (communiste), par le Vietnam, chassant ainsi les khmers rouges.
		Ceci discrédite alors la thèse selon laquelle les proximités idéologiques empêchent tout conflit, ce qui aurait rendu obsolète la géopolitique ici.
		Les chinois lancent alors une offensive dans le Vietnam du Nord le mois suivant.

\section{Les concepts de la géopolitique contemporaine}

	\subsection{Le relativité de l'influence des milieux naturels sur les orientations, les rivalités et les enjeux politiques}

		\paragraph{Une nouvelle approche}

		La géopolitique étudie les interactions entre l'espace géographique et les rivalités politiques dont il est le théâtre et souvent l'enjeu.

		À ces rivalités, l'espace géographique offre des opportunités ou impose des contraintes.
		Mais les unes comme les autres ne sont pas immuables.
		Leur influence dépend des capacités (démographiques, technologiques, financières) des pouvoirs qui convoitent ou défendent cet espace.

		Les rivalités ne sont pas seulement fondées sur des données objectives.
		En retour, ces rivalités influencent l'espace géographique.

		\paragraph{Opportunités}

		Exemples de positions stratégiques :
		\begin{itemize}
		\item Homs, sur une route marchande historique,
		\item Briançon, sur un verrou glaciaire, au carrefour de cinq vallées, citadelle fortifiée par Vauban,
		\item Lyon, anciennement Lugdunum, carrefour aussi, avec plusieurs accès fluviaux, d'où son importance à l'époque romaine.
		\end{itemize}

		Les positions stratégiques ne sont pas toutes situées en hauteur.
		Les gués (\textit{ford} en anglais) en sont aussi.
		Les fleuves également qui sont propices au transport des marchandises et constituent les berceaux des civilisations, le Nil et l'Euphrate notamment.

		On a aussi des contraintes à prendre en compte.
		Exemples :
		\begin{itemize}
		\item Accessibilité : Ardennes infranchissables selon Jules César, opinion répandue jusqu'à ce que les allemands les utilise comme point de passage pendant le seconde guerre mondiale,
		\item Climat froid : campagne de Napoléon en Russie, retour épouvantable en raison des conditions climatiques. Allemands après la défaite de Stalingrad dans des conditions similaires.
		\item Déserts : Sahara, Taklamakan (en Chine).
		\item Contraintes hydrographiques : inondations, crues du fleuve jaune en Chine.
		\end{itemize}

		Mais les contraites et les opportunités ne sont pas immuables.
		Des contraintes ont ainsi pu être surmonté avec la modernisation.

		De nouvelles contraintes apparaissent : accès à l'eau, contrôle des puits, usines de désalanisation.\\

		Tout dépend des capacités technologiques du moment.

		Exemples :
		\begin{itemize}
		\item Capacité à creuser des tunnels.
				On est passé en un siècle de la pioche au tunnelier, en passant par les marteaux piqueurs.
		\item Construction des ponts qui donnaient aux romains un avantage géopolitique durant l'Antiquité.\\
		\end{itemize}

		Le facteur crucial est aussi celui des moyens financiers disponibles.

		Exemples :
		\begin{itemize}
		\item Canal de Néron en Grèce : perce l'isthme de Corinthe pour raccourcir le trajet des navires de plusieurs centaines de kilomètres, initié par Néron.
			Mais Néron mourra peu après et le projet est stoppé car l'empire romain n'a pas les moyens de poursuivre un aussi gros chantier.
			Le canal ne sera percé que longtemps plus tard par des français, avec des capacités technologiques plus importantes.
		\item La Grande Muraille de Chine (début de le construction sous la dynastie Qin) :
			la Chine impériale a les moyens de mobiliser des fonds et des forces comme l'empire romain ne le pouvais pas.
		\item Canal de Suez : coût de près de quatre milliards d'euros, réalisation impossible avant la révolution industrielle.\\
		\end{itemize}

		Elles dépendent aussi de la volonté politique de cette mise en œuvre.

		Exemples :
		\begin{itemize}
		\item Pont-tunnel entre la Suède et le Danemark : les deux pays possèdent à la fois l'argent, la technologie et les compétences.
			Le chantier aurait donc pu être effectué depuis longtemps.
			Or il n'a été réalisé qu'au début de ce siècle car la Suède possède une forte culture insulaire.
		\item Le tunnel sous la manche : très violente campagne en Angleterre contre le tunnel, notamment de la part des milieux nationalistes lors de la proposition initiale du projet.
			Le projet ne démarrera que dans les années 80, avec des moyens technologiques encore accrus.
		\item Ligne ferroviaire Beyrouth-Haïfa : tunnel de Rosh Ha Nikra.
			Le tunnel est clos aujourd'hui car toute circulation entre le Liban et Israël est prohibée.
		\end{itemize}

		\begin{quote}
			« On ne voit que ce qu'on regarde et on ne regarde que ce que l'on a déjà dans l'esprit. »
		\end{quote}

		L'imaginaire individuel et social est fasciné par les lieux.

	\subsection{Les représentations : les lieux et représentations mystiques et symboliques}

		\paragraph{La construction des représentations}

		Pour certains groupes on trouve des croyance ou des mythes fondateurs.
		Certains sont d'ailleurs revendiqués par plusieurs groupes.

		Exemple : lieux revendiqués à la fois par le judaïsme et l'islam, comme le cénotaphe de Jacob ou le dôme du rocher.

		Autre mythe fondateur important : l'empereur jaune en Chine.
		L'Empereur jaune, Huangdi, qui aurait régné de 2697 à 2597 av. J.-C. est considéré comme le père de la civilisation chinoise, particulièrement à partir du XIXe siècle où la définition de la nation chinoise fait l'objet de nombreux débats.

		\paragraph{La construction d'un mythe national : le cas de la France.}

		On trouve différentes strates majeures dans l'élaboration d'un mythe national français :
		\begin{enumerate}
		\item les Gaulois,
		\item Charlemagne (mais revendiqué aussi par les allemands),
		\item Jeanne d'Arc,
		\item la révolution,
		\item Napoléon,
		\item première guerre mondiale,
		\item la résistance et la victoire dans la seconde guerre mondiale.
		\end{enumerate}

		L'histoire est enseignée de sorte à constituer un corpus qui va forger une représentation.
		La synthèse de cette histoire nationale va se retrouver dans les symboles : drapeau, devise, l'hymne, les commémorations (fête nationale,...).

		\paragraph{Construction et déconstruction de l'identité nationale.}

		Les États jouent sur l'altérité : se construire une identité qui nous permette de nous différencier des voisins, une altérité trop forte conduisant au conflit.

		À la base de l'identité on retrouve la cellule familiale avec transmission d'une identité et les structures claniques.
		Le rôle de l'école est également clé dans la construction d'une identité à un niveau national, avec un biais souvent dans les livres d'histoire.

		Le sentiment national est renforcé par l'armée : service militaire, parades, glorification par les médias nationaux.

		\paragraph{Le rôle des dénominations.}

		On trouve également un ancrage identitaire dans les dénominations et la cartographie.
		Exemple : nom des villes aussi indiqués en breton ou en alsacien.

		Des actions peuvent d'ailleurs être intenté contre l'identité même des villes.
		Exemples :
		\begin{itemize}
		\item révolution qui voulait raser Lyon et le renommer “Ville affranchie”,
		\item Jérusalem interdite aux juifs par les romains et renommée Ilia,
		\item noms des villes changés sous la Russie bolchevique : de St-Petersbourg à Leningrad.
		\end{itemize}

		\paragraph{Cartes réelles et cartes mentales.}

		Les cartes étaient traditionnellement élaborées par l'armée.
		On trouve sur certaines cartes anciennes des indications pour les artilleurs.

		Aujourd'hui encore certains États rechignent à laisser accès au grand public aux cartes de certaines régions.

		Des cartes imaginaires peuvent aussi rentrer dans les mœurs, qu'elles soient purement fictives (carte de l'île au trésor) ou basées sur des ambitions politiques (carte d'une Allemagne élargie).

		Les cartes mentales et cartes réelles peuvent être très différentes dans la mesure où les cartes mentales reposent sur des présupposés qui occultent des points clés de la réalité.

		Ces représentations touchent évidemment aux régions les plus sensibles du monde, telle la Palestine.

		\paragraph{Les représentations « icônes » et repoussoirs.}

		On trouve d'abord la représentation du pays « à soi », qui se trouve souvent personnifiée (Marianne, oncle Sam, anglais opulent).
		Le symbole peut être autre : hexagone pour la France, botte pour l'Italie.

		Viennent ensuite les pays amis, avec là encore des représentations symboliques de référence.
		Par exemple la statue de la liberté, cadeau de la France aux États-Unis.

		On trouve en face premièrement les pays menaçants.
		Parmi la symbolique associée on trouve par exemple celle du bolchevik avec le couteau entre les dents.

		On trouve enfin les pays ennemis, historiquement le Royaume-Unis pour la France, puis l'Allemagne.
		Aujourd'hui l'ennemi est devenu Daech.

		\paragraph{Les représentations, elles non plus, ne sont pas immuables.}

		Ces représentations évoluent au gré des alliances.
		L'Allemagne, hier ennemie de la France est ainsi vue aujourd'hui comme sa meilleure amie.

	\subsection{La sédimentation des politisations successives de l'espace.}

		Exemple du sens des trains dans l'Est de la France, possédée un temps par l'Allemagne : l'ancienne frontière est encore aujourd'hui visible par le changement de côté du train.
		On trouve sur le même modèle l'exemple du concordat en Alsace-Lorraine.

		Cette sédimentation peut se faire sur des régions façonnées pendant longtemps par une culture avant de passer à une autre.
		C'est le cas de l'Andalousie, qui a été pendant longtemps musulmane, avant que le dernier calife ne soit chassé par les rois catholiques (fin de la conquête avec la prise de Grenade).
		Des monuments comme la cathédrale de Séville ou la grande mosquée de Cordou présentent ainsi un profil mixte.

		Dans l'autre sens on trouve l'église de Sainte-Sophie, érigée à Constantinople, qui fût premièrement une église (construite au VI\up{ème} siècle), puis a été transformée en mosquée et enfin en musée par Ataturk dans une perspective laïque.
		Un retour au statut de mosquée serait même envisagé par le gouvernement actuel.

		On trouve enfin l'exemple de l'esplanade des mosquées, ou Mont du Temple.


\section{Les « objets » de la géopolitique : territoires, frontières et flux.}

	\subsection{En terme de géopolitique, le territoire national c'est l'espace de l'État.}

		\paragraph{Les origines de l'État.}

		Même dans les sociétés sans État, on trouve des problèmes de territoires.
		On peut penser notamment à certaines tribus\footnote{
			Ethnies : présentent une homogénéité culturelle, du domaine de l'idéel, exemples : iroquois, mohicans.
			Tribu : réseau de solidarité qui va donner des droits de chasses et la défense d'un territoire.
			}
		en Amazonie ou autres territoires peu construits.
		Avant la révolution néolithique étaient pratiquées notamment la chasse, la cueillette et la pêche, le tout de façon nomade.
		Va ensuite se développer l'agriculture et la domestication du bétail.
		Le prix à payer est que le travail du paysan est plus important que celui du chasseur.

		Ces paysans vont alors stocker des denrées (céréales,...), ce qui crée donc le besoin de protéger ces stocks.
		Il faut donc des soldats, mais également une autorité sur ces soldats.
		Des ressources comme l'eau vont exiger également la présence d'une autorité.
		Cette autorité sera de plus connectée au ciel car c'est lui qui influe sur la production agricole.

		La révolution néolithique va ainsi engendrer la création de souverains qui sont en même temps grands prêtres.
		Ce souverain est donc à la tête d'une armée, d'une religion et d'une civilisation.

		S'élaborent alors petit à petit des codes qui régissent le territoire.

		\paragraph{L'objectif de l'État est d'assurer sa pérennité et d'affirmer sa légitimité.}

		La légitimité doit être affirmée à la fois vis-à-vis du territoire, et celle de sa population.
		Ceci se fait également à travers l'affirmation d'une altérité.

		\paragraph{Pour assurer cette pérennité, l'État se donne pour mission de : protéger le pouvoir, protéger le territoire national, séduire, surveiller, punir.}

		L'exercice de l'État se fait souvent dans des lieux protégés (Cité interdite, Kremlin).

		La protection du territoire se fait à travers des aménagements (Vauban par exemple sous Louis XIV), qui peut prendre le pas sur les objectifs de conquête.

		L'État doit également valoriser le territoire.
		Exemple : Noordoostpolder aux Pays-Bas, territoire repris sur la même.
		Ceci se fait également par la valorisation des compétences et de la production.

		L'État déploie également des instruments de séduction : organisations d'événements (sportifs, célébrationnels,...), interventions extérieures à image positive (humanitaire, partage de progrès technologique), investissements.
		Il peut chercher particulièrement à impressionner, par sa richesse : importance des palais (Versailles, Westminster), souverain, ou par sa puissance militaire (défilés).

		L'État cherche également à surveiller, à la fois à l'intérieur et à l'extérieur en espionnant.

		Enfin l'État doit punir, avec notamment pour les individus la peine de mort ou la réclusion (bagne de Cayenne).
		On trouve également des punitions entre États : intervention militaire, arme atomique.

		Pour cela l'État doit financer sa puissance, par les impôts, les prises de guerre ou les accords.
		Même dans les guerres modernes, le vaincu doit faire des promesses d'indemnité envers le vainqueur (traité de Versailles, pacte de Quincy en 1945 en les États-Unis et la dynastie Saoud).

	\subsection{La frontière est une discontinuité géopolitique, à fonction de marquage réel, symbolique et imaginaire}

		\paragraph{Une discontinuité qui évolue avec le temps.}

		Dans le cas de marquage géographique, des évolutions du terrain pour des causes naturelles ou artificielles menacent de modifier les frontières.

		\paragraph{Des marches à la frontière linéaire.}

		La frontière n'est pas linéaire initialement, il existe des enclaves et autres particularités géographiques.
		Il existe l'expression d' « esprit de frontière » pour décrire une vision mentale idéalisée de la frontière.
		Beaucoup, comme Danton, sont convaincus que ce sont des frontières naturelles qui délimitent le territoire français : Rhin, chaîne des crêtes dans les massifs montagneux.

		Ceci peut donner des situations particulières : cas de Saint Gingolph.

		\paragraph{Fonction de marquage réel.}

		La fonction de marquage réelle peut très bien être peu visible dans le cas de frontières ouvertes : bornes discrètes.

		Elle peut aussi être filtrante, avec différents niveaux :
		\begin{itemize}
		\item petit poste frontière à Saint-Gingolph, peu filtrant,
		\item poste de frontière de San Ysifro entre l'Amérique et le Mexique, réputé être le plus fréquenté du monde, contrôlé pour maîtriser les migrations,
		\item frontière entre l'Inde et le Pakistan : un seul poste à Wagah, qui n'est ouvert que dans la journée.
		\end{itemize}

		On a enfin les frontières fermées, qui peuvent être des parties de frontières filtrantes :
		\begin{itemize}
		\item entre les deux Corées,
		\item entre Israel et le Liban,
		\item ligne de cessez-le-feu dans le Golan entre Israel et la Syrie,
		\item entre l'Inde et le Pakistan.
		\end{itemize}

		Pour certaines frontières comme entre l'Algérie et le Maroc il y a danger mortel à toute traversée.
		Les fermetures de frontières ne sont pas forcément définitives ni de très longues durée : exemple de fermetures temporaires pendant l'épidémie Ebola entre les pays concernés.

		\paragraph{Fonction de marquage symbolique}

		On trouve un marquage symbolique dont l'utilité est de faire comprendre de quel côté de la frontière l'on se situe : présence du drapeau, uniformes des douaniers différents entre pays frontaliers.

		\paragraph{Fonction de marquage imaginaire : la frontière vécue}

		La frontière peut-être vécue d'une façon particulière :
		\begin{itemize}
		\item chemin des toblerones : défenses anti-char de la Suisse,
		\item le château de Voltaire, sur la frontière franco-suisse, pour pouvoir fuir le pays facilement,
		\item camps de réfugiés syriens.
		\end{itemize}

		Cette frontière vécue peut être frustrante :
		\begin{itemize}
		\item même tribus des deux côtés de la frontière entre le Maroc et l'Algérie,
		\item clôtures et rideau de fer, à Berlin ou à Tel-Aviv,
		\item frontières dissymétriques, des niveaux de vie différents des deux côtés. Exemples : Tijuana entre le Mexique et l'Arizona, ou Melila, une enclave espagnole sur la côté marocaine.
		\end{itemize}

		L'absence de frontière peut faire peur : frontière “passoire”.
		Cela touche l'imaginaire de ceux qui en vivent loin.

		Certains rêvent de frontières qui n'existe pas, manifestant un désir d'indépendance.
		Exemples :
		\begin{itemize}
		\item Québec,
		\item Catalogne,
		\item Kurdistan : aspiration des populations kurdes, état quasi-indépendant sur la partie irakienne.
		\end{itemize}

		On a également l'irrédentisme, le désir d'extensions des frontières.
		Exemple : grande Hongrie, courant nationaliste haineux envers les français à cause du traité de Trianon de 1920.

		\paragraph{Les frontières ne sont pas fixées par Mère Nature}

		Les frontières sont la ligne d'équilibre à un moment donné entre États voisins (frontière franco-espagnole par exemple) ou entre puissances (frontières du Proche-Orient ou du continent africain, résultats du rapport de force entre le Royaume-Uni et la France).
		Tous les États d'Europe de l'Est ont également des frontières héritées de rapports entre puissances historiques.

		Mais, bien entendu, les frontières suivent, à chaque fois que cela est possible, une ligne naturelle de défense (rivière ou ligne de crêtes par exemple).

		Dans la façon dont s'étendent les frontières, on passe ainsi souvent d'une frontière naturelle à une autre.

		Exemple de la Géorgie : territoire des ossètes dont les désirs de rejoindre la Russie sont entretenus par celle-ci.

		Cas de l'Ukraine : tout l'Est est russophone, de même que la Crimée.
		Le berceau du peuplement russe est l'Ukraine, c'est pourquoi les russes sont particulièrement sensibles à sa perte.

		En Asie centrale ont lieu des contestations au niveau des frontières dessinées par l'URSS : Kazakhstan, Turkménistan, Ouzbékistan, Tadjikistan, Kirghizstan.
		Plusieurs frontières y sont discutées.

		\paragraph{Plusieurs types d'évolution des frontières possibles}

		\begin{itemize}
		\item La scission : déroulement assez pacifique.
			Exemples : Norvège qui se détache de la Suède en 1905, Tchécoslovaquie divisée lors de la dislocation de l'URSS.
		\item La partition : déroulement assez violent, nommée aussi sécession par ceux qui ne la souhaitent pas.
			Exemples :
			\begin{itemize}
			\item[\textbullet] partition des Indes britanniques en Inde, Pakistan et Pakistan Oriental (devenu Bangladesh en 1957 à la suite d'une seconde partition),
			\item[\textbullet] partition de l'Irlande avec la province d'Ulster pour des raisons religieuses : population à majorité protestante au Nord e catholique au Sud,
			\item[\textbullet] Yougoslavie, décidée à Versailles, qui s'est disloquée en 1991,
			\item[\textbullet] Corée en 1953.
			\end{itemize}
		\item L'annexion : différentes façons possibles.
			\begin{itemize}
			\item[\textbullet] par la force : Sahara occidental par le Maroc en 1975, ou Crimée par la Russie en 2014,
			\item[\textbullet] par achat : Alaska en 1867 achetée à la Russie, ou Louisiane en 1803 achetée à la France (territoire plus large que l'État de Louisiane actuel),
			\item[\textbullet] par référendum, après accord : Savoie et Nice en 1860.
			\end{itemize}
		\end{itemize}

	\subsection{Les flux : humains, économiques, immatériels}

		\paragraph{Les flux humains.}

		Ces flux posent des problèmes géopolitiques majeurs dans les pays de départ, au zones de passages qui sont en tension et dans les pays d'arrivée.
		On trouve des flux internationaux et des flux régionaux.

		\paragraph{Les flux économiques / marchands.}

		Ils donnent à certains points de la planète une importance géostratégique majeure.
		Exemples : les détroits, comme celui de Malaka entre Singapour et l'Indonésie, le canal de Suez, le collier de perles (l'ensemble des ports où la Chine a pris position en Asie du Sud pour assurer son commerce).

		\paragraph{Les flux énergétiques.}

		Ces flux se font à travers des pipelines (oléoducs ou gazoducs).
		Par exemple la création d'un gazoduc sous la baltique par les russes pour desservir d'Allemagne court-circuite les pays baltes ou la Pologne qui sont entre les deux.

		\paragraph{Les flux d'informations.}

		Exemple : espionnage par les américains grâce à leur maîtrise des technologies actuelles.




\part{Principaux aspects de la géopolitique du monde actuel}

\section{Un constat : mondialisation versus fragmentation politique}

	\subsection{Mondialisation}

		\paragraph{Les évolutions technologiques et leurs conséquences géopolitiques.}

			La mondialisation est une conséquence de la révolution industrielle, et en particulier dans les transports avec la machine à vapeur qui permet la locomotive et l'extension du chemin de fer.
			Avant cela tout voyage était pénible et dangereux, en plus d'être cher.
			Avec le chemin de fer le voyage devient un plaisir, et plus abordable.
			Commencent alors à se poser des problèmes de souverainetés : passeport, etc...

			Le chemin de fer est alors au cœur d'enjeux stratégiques : réseau ferré très centralisé et qui répond aux besoins des militaires.
			La Russie choisit l'écartement de ses rails large pour des raisons stratégiques.

			On assiste au début de l'automobile à la fin du XIX\up{e} siècle.
			La première Panhard et Levasser en 1890 roule à 20 km/h.
			Au début du XX\up{e} siècle Bugatti sort un modèle qui atteint les 200 km/h.
			Le réseau des autoroutes est également centralisé autour de Paris, surtout dans un premier temps.

			Vient également le développement de l'aviation.
			Les frères Wright effectuent leur premier vol motorisé en 1903 et Louis Blériot réalise le première traversée de la Manche en 1909.
			Les dirigeables se développent également, surtout du côté allemand avec l'entreprise Zeppelin.
			Se pose alors la question de savoir à qui appartient l'espace aérien.

			La révolution des télécommunications révolutionne également le transport des idées.
			Au télégraphe optique de Chappe succède la télégraphe électrique.
			Ce développement se fait tout d'abord au service de l'État.

			Puis vient le téléphone ainsi qu'une amélioration du télégraphe par Belin.
			Arrive ensuite la radio-transmission avec la radio, puis la télévision qui est tout d'abord un média d'État.
			Tous les grands pays ont alors des chaînes d'ondes courtes, capables de se propager à toute la planète.

			On commence ensuite à exploiter l'espace : premier satellite de télécommunication, le Telstar, envoyé en 1962 (Spoutnik en 1957).
			Problème là encore pour les nations : à qui appartient l'espace.
			Est donc établit le traité de l'espace en 1967.

			Aujourd'hui le monde est connecté grâce à internet (un peu plus de la moitié de la population mondiale).
			L'invention a été faite au CERN puis le développement aux USA, qui s'est octroyé depuis le début des pouvoirs importants sur les instances de l'internet.
			Ils sont également à l'origine de nouvelles formes de communication : messagerie, réseaux sociaux...

			Ce manque de contrôle des autres pays sur internet et ces moyens de communication pose des problèmes de souveraineté.
			C'est pourquoi on trouve par exemple en Chine des alternatives à Facebook, Google ou Twitter.

		\paragraph{Centre et périphérie.}

			\subparagraph{Économiques}

			Le centre est usuellement assimilable à l'Occident mais l'appartenance à se centre se discute.
			Qu'en est-il par exemple de la Chine aujourd'hui ?

			Les grandes entreprises actuelles sont présentes sur toutes la planète.

			On a de plus une compétition acharnée entre pays du centre.
			Exemple : Airbus et Boeing, le second étant très soutenu par les USA.
			Les américains ont de plus tenté de saboter l'industrie spatiale européenne.

			La compétition se fait aussi sur le commandement économique : on essaie d'avoir chez soi les centres de commandement des grands groupes, qui vont de pair avec les emplois de plus hautes compétences, un attrait et une influence.

			À l'ère des empires coloniaux le centre, premièrement européen, profitait d'un avantage technologique et de la possession des colonies.
			Il dominait militairement et pouvait imposer ses produits sur le marché mondial, avec des conséquences négatives pour les producteurs locaux.

			Aujourd'hui, malgré la fin des empires coloniaux, le centre continue à exercer une pression sur la périphérie.
			Exemples :
			\begin{itemize}
			\item interventions des USA en Amérique latine, avec force assassinats et coups d'état pour garantir leur prédominance sur le marché, comme l'installation de la dictature de Pinochet au Chili, le plus gros producteur mondial de cuivre.
			\item industrie du textile qui fait travailler des ouvriers en Asie, le Bangladesh étant actuellement le moins cher,
			\item catastrophes humaines et écologiques : immeuble du Bona Plaza qui s'effondre ou usine de Union Carbide au Bhopal qui explose.
			\end{itemize}

			On trouve aussi d'autre périphéries plus particulières, comme les paradis fiscaux, dont les plus opaques se trouvent géographiquement dans la périphérie et sont les plus corrompus.
			Leur développement se fait avec les mafias internationales.

			\subparagraph{Culturels}

			La mondialisation est également culturelle.

			Exemple : implantation inégale de Mac Donalds en fonction des niveaux de vie locaux.

			L'anglais s'impose comme langue universelle.
			Cela suscite la réaction de d'autres pays qui tiennent à leur influence culturelle : institue Confucius, Goethe Institut, institut français...

			La diffusion culturelle est également mondiale : France 24, CCTV...

	\subsection{Fragmentation}

		On assiste à une multiplication des souverainetés.
		Il y avait 51 signataires de la charte de l'ONU en 1945, puis 100 dans les années 60, et aujourd'hui on compte 193 pays membres et 2 pays observateurs (Vatican et Palestine).

		\paragraph{Décolonisation.}

		Dans l'Afrique de l'après-guerre seul un état était indépendant, l'Éthiopie.
		Après 1965 on trouve tout une mosaïque d'États.

		\paragraph{Dislocation.}

		Se disloquent notamment les anciens pays communistes : URSS et Yougoslavie.

		\paragraph{Balkanisation et libanisation}

		Les Balkans sont des anciens pays de l'empire ottoman avec un découpage ethnique très complexe, sans continuité géographique des populations, voire mélangées.

		La balkanisation désigne ainsi la situation dans laquelle le populations aspirent à créer un État homogène avec le groupe auquel elles appartiennent.
		Cela amène à une situation de conflit permanent et, dans le cas des Balkans mêmes, la répartition de la population est trop complexe pour permettre de trouver une solution.

		La situation au Liban est celle d'un grand brassage de religions : catholiques, orthodoxes, musulmans sunnites, musulmans chiites, druzes...
		On parle alors de libanisation quand l'une des parties veut gouverner le tout.
		À la chute de l'empire ottoman la Société des Nations avait accordé à la France et au Royaume-Uni des mandats sur les territoires arabes qui lui appartenaient.
		En particulier le Liban a été créé par les français de sorte que les catholiques (marronites) puissent gouverner avec une majorité mais la démographie a changé depuis et, la majorité n'ayant plus cours, on a un partage du pouvoir.

		On peut rapprocher cela de différents cas :
		\begin{itemize}
		\item en Afrique : découpage qui ne tient pas compte de la répartition des tribus,
		\item Afghanistan : patchwork de populations iraniennes, turques ou autres, avec à l'intérieur de chaque population des tribus, dont les Pachtounes (la plus étendue, iranienne).
			On a une libanisation car les talibans ne revendiquent pas le contrôle de la seule zone Pachtoune mais de tout l'Afghanistan.
		\item Syrie : libanisation avec les alaouites (hérésie du chiisme) qui gouvernent, les kurdes (non arabes) au Nord-Est du pays, des musulmans sunnites, des druzes au Sud...
			Hafez El Assad a installé des chiites à tous les postes étatiques importants : armées, services secrets, administration...
			Le régime mis en place fait montre d'un contrôle absolu sur la population ou les communications, avec une grande violence.
		\item Côte d'Ivoire : présence de quatre grands groupes ethno-religieux.
		\end{itemize}

		Cela peut générer des déplacements de populations ou des phénomènes d'épuration ethnique, comme dans les balkans dans les années 1990.

		\paragraph{Séparatisme et indépendantisme}

		On distingue des deux phénomènes précédents les revendications séparatistes, ou indépendantistes (selon le point de vue), que l'on trouve dans de nombreuses régions en Europe par exemple.
		Les facteurs à l'origine de ces revendications peuvent être religieuses, mais aussi économiques (possession de matières premières qui fournissent un avantage sur le reste du pays par exemple).

	\subsection{De nouveaux acteur géopolitiques}

		\paragraph{Les ONG}
			ont acquis au cours des trente dernières années une légitimité sur la scène internationale, de sorte que même les pays qui semblent peu se soucier de leur image à l'international peuvent difficilement ignorer les campagnes de certaines d'entre elles.

		\paragraph{Les organismes supranationaux :}
			le conseil de sécurité de l'ONU (le “machin” selon Charles de Gaulle), ou la cour pénale internationale.

		\paragraph{L'OMC}

		\paragraph{Les réseaux mafieux}
			dont le succès se fait au détriment de la richesse des états.

		\paragraph{Les réseaux terroristes islamistes}
			au Moyen-Orient et en Afrique, qui usent des moyens technologiques actuels pour s'organiser.

		Les structures supra-nationales existent depuis longtemps (exemple : influence de l'Église catholique pendant des siècles).

\section{Instabilités du monde contemporain}

	\subsection{Démographiques : plus de 7 milliards d’êtres humains}

		La population mondiale est de 7 400 000 000 fin 2015.

		La Chine et l'Inde devrait prochainement atteindre chacun une population de l'ordre de $1{,}5$ milliards.
		D'ici à 2050 on devrait observer notamment une croissance démographique en Afrique.

		Ceci se traduit de façon inégale au niveau des matières premières.
		Les prix ne suivent pas forcément l'augmentation de la demande qui est observée.

		Des pays cherchent pour ses raisons à acquérir des terres arables pour subvenir à leurs besoins.
		Ainsi de nombreuses terres sont achetées en Afrique par l'Amérique du Nord, l'Europe, l'Asie, le Moyen-Orient, l'Afrique elle même, ou bien les autres.

	\subsection{Économiques : ressources limitées, besoins croissants}

		On avait en fait jusqu'en 1700 un PIB par habitant proche dans les différentes zones de la planètes.
		À la suite de la révolution industrielle, et notamment au cours du XX\up{e} siècle, ce PIB par habitant a évolué de façon très différente dans les différentes zones, celui de l'Amérique du Nors et de l'Europe de l'Ouest augmentant beaucoup plus vite que le reste.

	\subsection{Socio-culturelles : pertes et quêtes d’identité}

		Ceci s’exprime d'abord à travers la langue.
		Il existerait aujourd'hui entre 3000 et 7000 langues vivantes (selon les critères retenus).
		Une langue disparaîtrait toutes les deux semaines.

		On trouve aussi des problèmes identitaires.
		Exemple : différence culturelle importante en Papouasie-Nouvelle Guinée entre les populations papoue et guinéenne.

		\paragraph{Retour du religieux ?}

		Alors que la religion était largement attaquée dans les années 60 par exemple, ce phénomène s'est largement atténué aujourd'hui, et la religion trouve ses défenseurs.
		On pourrait donc évoquer un retour du religieux aujourd'hui.

	\subsection{Environnementales : réchauffement climatique, pollution, biodiversité menacée, etc.}

		Certaines régions ont toujours été plus à risque que d'autres : risques sismiques, volcans...
		Sont venues s'ajouter à cela des responsabilités dites anthropiques, dues notamment au rejet par l'homme de gaz à effet de serre.

		Cela menace la biodiversité de la planète.
		La pollution se retrouve dans les océans avec la constitution de vastes zones de déchets, ce qui gêne le développement du phyto-plancton.
		Enfin on trouve notamment l'élévation de la température mondiale, qui pose un problème géopolitique majeur.
		Des populations sont notamment amenées à se déplacer en raison de la hausse du niveau de la mer qui y est liée.

	\subsection{Politiques : un monde unipolaire, bipolaire, multipolaire ?}

		\paragraph{La fin de l’hyperpuissance américaine ?}

		Les États-Unis, qui jouissaient du statut d'hyperpuissance, voient leur statut remis en question après des échecs comme l'intervention en Irak.
		Leur PIB, qui vient d'être rattrapé par la Chine, continue toutefois de croître, tout en étant déjà très haut.
		Les États-Unis gardent de plus un avantage important en ce qu'ils concentrent un grand nombre de firmes majeures, les meilleures universités et centres de recherche (beaucoup de prix Nobel).

		Un marqueur important de la puissance américains est le poids du dollar dans les systèmes internationaux.
		Les accords de Bretton Woods (New Hampshire) furent signés le 22 juillet 1944, établissant la supériorité du dollar.
		À la suite de la guerre les États-Unis possédaient une large partie des ressources en métaux précieux, créant un déséquilibre avec le reste du monde.
		Est donc créé l'étalon de change or : l'once d'or vaut 35\$ et ce taux est fixe.
		En effet il n'y plus assez d'or dans les banques centrales pour suivre l'augmentation du besoin de monnaie, d'où le passage par le dollar.

		Enfin les États-Unis sont une puissance militaire : près de 37\% des dépenses militaires de la planète.
		La Chine arrive deuxième avec 11\%.
		Ils ont notamment l'avantage sur leurs concurrents en matière d'avancée technologique.
		Là où la Chine possède une armée classique, les États-Unis investissent massivement dans des technologies de pointe : drones, robots...
		L'US Navy, la marine de guerre de États-Unis, est la première force naval au monde, donc la portée et la force de frappe excède de loin celle des autres pays.
		Ainsi ils possèdent en permanence dix porte-avions en mer, la deuxième puissance en la matière étant la France avec un porte-avion à propulsion nucléaire.

		Les américains manient de plus le soft power : média, industrie du divertissement, outils de communications, livres...

		\paragraph{La Chine, pays émergent ?}

		Puissance militaire.

		Puissance spatiale.

		Grande disparité : encore des zones avec des équipements très faibles.
		Xinjiang avec un statut particulier : mouvement sécessionniste durement réprimé, population Ouighur, région qui progresse économiquement en raison de l'immigration Han.

		\paragraph{Les autres BRICS}

		Ville de Bengalore : présente des contrastes importants.

		Brésil.

		Russie : vit principalement de la vente de ses matières premières.
		Pauvreté parfois très marquée.

		Questionnements sur l'avenir de l'Afrique : région du monde avec les plus forts taux de croissance.
		La situation est la aussi contrastée.

		\paragraph{Demain, l’Afrique ?}

\section{Les foyers d'instabilité}

		Les foyers en question peuvent être relatifs à des instabilités à plus ou moins long terme.

		\begin{itemize}
		\item L'Afghanistan. De nombreux pays ont participé une opération visant à éliminer les talibans.
		\item L'Afrique orientale
		\item Le Moyen-Orient
		\end{itemize}

\section{Le principal foyer d'instabilité dans le monde : le Moyen-Orient}

	\section{Proche-Orient, Moyen-Orient}

		Le Moyen-Orient s'étend de l'Égypte jusqu'à l'Iran.

		Chez les américains le Moyen-Orient comprend également le Pakistan et l'Afghanistan.
		Ce que Bush avait appelé le grand moyen Orient s'étend même jusqu'au Maroc à l'Ouest.

		On distingue également en français le Proche-Orient ainsi que le levant.
		Le Proche-Orient comprend l'Égypte, la Jordanie, l'Irak, la Syrie, le Liban et l'ensemble constitué d'Israël et des territoires palestiniens.
		C'est la partie arabe de l'empire ottoman, puisque ces territoires étaient sous la souveraineté du sultan de Constantinople jusqu'à la première guerre mondiale.

		La région ne peut pas être caractérisée par une unité géographique ou climatique particulière : présence de reliefs variés, climat globalement aride mais avec des zones comme le Liban qui so
		D'autres pays bénéficient grâce à des fleuves « l'eau que le ciel leur refuse », comme l'Égypte avec le Nil.

		Il n'y a pas non plus, et c'est la clé du problème, d'unité humaine, ce depuis des millénaires.
		On trouve des populations arabes, perses, turques, juives et kurdes.
		Tois monothéismes sont présents mais la religion musulmane est majoritaire.

		Se superposent une complexité ethnographique et une complexité religieuse.
		En effet, au Moyen-Orient, le chiisme représenté une population importante alors que le sunnisme est largement majoritaire dans les autres régions musulmanes de la planète.
		Les chiisme sont en particuliers majoritaires en Iran.
		On trouve également beaucoup de communautés minoritaires issues du chiisme comme les druzes ou les alaouites.

		Là où les violences sont les plus fortes c'est donc au sein d'une religion donné, de même que les guerres de religion en France se sont faites entre chrétiens.

		Au Moyen-Orient on trouve également une population chrétienne, en particulier au Liban où ils constituent 40\% de la population.
		Mais cette population chrétienne est également divisée et certaines Églises sont dites autocéphales, i.e. indépendantes des Églises principales, comme les coptes en Égypte.

		On trouvait également de grands foyers de population juive, comme à Baghdad, à Sanaa,...
		On a alors assisté à un retour d'une communauté juive à partir de la fin du XIX\up{e} siècle, puis avec la création d'Israël.
		Aujourd'hui on trouve une forte population juive au Moyen-Orient mais concentrée en Israël car les juifs présents dans d'autre pays arabes ont été amenés à émigrer plus ou moins volontairement.

	\subsection{Enjeux stratégiques et conflits à portée mondiale}

		Cette région a toujours été à la fois au cœur d'enjeux majeurs et le théâtre d'affrontement de grandes puissances du monde.

		Dans les temps anciens cette région était convoitée pour des raisons commerciales : position stratégique et production de biens précieux.

		À la fin de l'Antiquité cette région commence à intéresser des empires plus éloignées : héllenisation du Moyen-Orient par les conquêtes d'Alexandre le Grand, puis appartenance à l'Empire romain.

		Le Moyen-Orient a pris au XIX\up{e} siècle une importance mondiale car les puissances européennes, Royaume-Uni et France en particulier, le voit comme un point stratégique.

		La pénétration européenne qui s'effectuait à l'époque de l'empire ottoman s'est affermit à la suite de la chute de celui-ci.

		Parmi les enjeux principaux aujourd'hui on trouve notamment le caractère de point de passage commercial (cana de Suez notamment) et la forte concentration d'hydrocarbure, celui-ci y étant particulièrement facile à construire.
		Pour des puissances comme les États-Unis il faut faire en sorte que le pétrole produit au Moyen-Orient ne vienne pas concurrencer celui produit sur le territoire américain.
		L'enjeu pétrolier est donc complexe et donc complexe car il intéresse le monde entier, et plus seulement les États-Unis et l'Europe, et intérêt se faisant à double niveau, avec un enjeu de prix.

	\subsection{Des conflits emblématiques}

		\paragraph{Le conflit israelo-palestinien}

		Ce conflit peut être qualifié de surreprésenté, ou de basse intensité.
		En effet le territoire concerné est de la taille de la Sardaigne, ou de la moitié du Danemark, et la population est de 10 millions environ, soit moins que l'agglomération parisienne.
		Enfin la totalité du conflit a fait moins de mort depuis 1947 que la seule journée du 21 août 1914 en France (bataille de Charleroy), ou que les narco-trafiquants au Mexique.
		Pourquoi ce conflit provoque-t-il donc autant d'intérêt et d'émotion dans le monde entier ?

		Ce conflit est en fait emblématique de beaucoup d'autre car c'est celui entre un monde développé et un monde en développement.
		La ligne verte définie en 1967 derrière laquelle est censée se tenir Israël est la « frontière de toutes les frontières ».
		En effet Israël possède un PIB par tête supérieur à celui de l'Italie et un IDH supérieur à celui de la France.

		La frontière est donc tout d'abord économique mais aussi culturelle.
		On a d'un côté une société occidentale fondée sur une hiérarchie sociale.
		Dans les territoires palestiniens la structure clanique pèse également dans les processus de décision.

		Les cessez-le-feu vont conduire à des lignes, dont la ligne verte, qui marquent la séparation avec les territoires arabes.


\end{document}